\documentclass[11pt,a4paper]{article}
\usepackage[T1]{fontenc}
\usepackage[utf8]{inputenc}
\usepackage{mathptmx}
\usepackage[scaled=0.92]{helvet}
\usepackage[a4paper,margin=2.4cm]{geometry}
\usepackage{booktabs,tabularx,array,enumitem,microtype}
\usepackage[hidelinks]{hyperref}
\usepackage{titlesec}
\titleformat{\section}{\sffamily\large\bfseries}{\thesection}{0.8em}{}
\titleformat{\subsection}{\sffamily\normalsize\bfseries}{\thesubsection}{0.8em}{}
\setlist{itemsep=2pt,topsep=4pt}
\setlength{\parskip}{0.45em}
\setlength{\parindent}{0pt}
\newcolumntype{Y}{>{\raggedright\arraybackslash}X}

\begin{document}

{\sffamily
{\small Aalto University, School of Science\\
Master's Programme in Computer, Communication and Information Sciences}\par\vspace{1.1em}
{\Large\bfseries Master's thesis plan}\par\vspace{0.5em}
{\large Scope Attribution for Static UI Transition Graph Extraction\\ in Single-Page Applications}\par\vspace{0.8em}
\begin{tabular}{@{}ll}
Student & Duy Khanh Tran, 103503394 \\
Supervisor & Prof.\ Petri Vuorimaa \\
Major & Computer Science \\
Extent & 30 ECTS \\
Language & English \\
Date & 22 September 2026 \\
\end{tabular}
}
\par\vspace{0.4em}\hrule\vspace{0.7em}

\section{Topic}

A UI transition graph models an application as a set of screens and the transitions
a user can trigger between them. Static extraction of such graphs from source code is
well established for Android but has not been studied for client-side web frameworks,
which are modelled dynamically by crawling. The thesis measures static extraction on
real Angular applications using an existing open-source analyser developed by the
student, classifies where it fails, and proposes and evaluates an analysis, called
scope attribution, that addresses the largest class of failures: navigation sites
located in shared components, injected services and route guards, which the current
extractor attributes to no screen.

\section{Research questions}

\begin{itemize}[leftmargin=2.4em]
  \item[\textbf{RQ1}] How much of the navigation in real Angular applications can be
        extracted statically? Measured as the fraction of navigation constructs in the
        source that become deterministic transitions in the graph.
  \item[\textbf{RQ2}] Where does extraction fail, and why? Answered by a defect
        classification in which each class must be confirmed by at least two
        independent evaluation methods.
  \item[\textbf{RQ3}] Does scope attribution recover the largest defect class? Measured
        as the change in the RQ1 ratio between two fixed tool versions on identical
        inputs, with precision on the deterministic tier not decreasing.
  \item[\textbf{RQ4}] Does the resulting graph detect source changes well enough to
        support change impact analysis? Measured by mutation testing.
\end{itemize}

The thesis does not claim completeness, which is undecidable for dynamic navigation;
it proves conditional soundness over a declared scope and measures recall. It makes no
claim about the effect of the graph on coding agents, and no claim beyond Angular.

\section{Material and methods}

\textbf{Corpus.} Six open-source Angular applications, selected by six published
criteria from ten candidates, pinned to exact commits, and spanning 14 to 674 declared
routes. Threats to the validity of an open-source corpus and their mitigations are
addressed in the thesis. An industrial case study on a proprietary Angular
application, reporting aggregate figures only, is proposed as an extension subject to
the supervisor's view and the employer's permission.

\textbf{Evaluation.} Four independent methods, each with a declared blind spot:
verification of every transition against its source witness (precision); comparison
with an independently written route parser (diagnosis); the applications' own passing
end-to-end tests as ground truth (recall); and mutation testing (sensitivity). The
whole evaluation is reproducible from one command with full provenance.

\textbf{Preliminary results.} On the current extractor, static extraction recovers
between 0 and 14 percent of navigation constructs as deterministic transitions, and
three of six applications yield routes but no transitions. Where deterministic
transitions are produced they are almost always correct (75 of 76 verified against
source). Three independent methods point to attribution as the dominant cause of loss.

\section{Planned structure}

\begin{center}\small
\begin{tabularx}{\textwidth}{@{}lYr@{}}
\toprule
 & Chapter & Pages \\
\midrule
1 & Introduction: motivation, research questions and scope, thesis structure & 4--6 \\
2 & Background: transition graphs; static extraction for Android; navigation models for
    web applications; dynamic models of single-page applications; Angular routing and
    dependency injection; formal foundations & 15--20 \\
3 & Research material and methods: corpus and validity; graph model; scope attribution;
    implementation; evaluation protocol & 15--20 \\
4 & Results, one section per research question & 12--15 \\
5 & Discussion: findings, limitations, declared approximations, future work & 5--8 \\
6 & Conclusion & 1--2 \\
\bottomrule
\end{tabularx}
\end{center}
Target length 60 to 75 pages. Chapters 1 to 3 are drafted.

\section{Schedule}

\begin{center}\small
\begin{tabularx}{\textwidth}{@{}llY@{}}
\toprule
Date & Milestone & Completion criterion \\
\midrule
28 Sep & Baseline frozen & Tool version tagged; corpus re-measured; result tables generated. \\
5 Oct & Chapter 3 final & Corpus, model and protocol sections complete. \\
12 Oct & Attribution implemented & Component nesting, router outlets, DI-aware call graph, guards, anchored edges. \\
15 Oct & Checkpoint & If attribution is incomplete, RQ4 moves to future work and RQ3 reports the implemented subset. \\
19 Oct & Second measurement & Same corpus and harness on the new tag; before and after tables. \\
26 Oct & Chapters 2, 4, 5 drafted & \\
2 Nov & Complete draft to supervisor & \\
November & Thesis presentation; revisions & Date agreed with supervisor. \\
30 Nov & Submission & PDF/A and abstract in MyStudies; graduation request in Sisu. \\
\bottomrule
\end{tabularx}
\end{center}

The student's study right ends on 31 December 2026. The schedule therefore targets
the December meeting of the Degree Programme Committee. The study coordinator will be
consulted on the exact date the study right must cover and on the extension
procedure.

\section{Risks}

\begin{center}\small
\begin{tabularx}{\textwidth}{@{}YY@{}}
\toprule
Risk & Response \\
\midrule
Attribution not complete by 12 October & The empirical study (RQ1, RQ2) stands alone. RQ4 is dropped first; RQ3 reports what is implemented and declares the rest. \\
Supervisor reading time in November & Draft date and reading window agreed in advance. \\
Corpus applications require a backend to run & No planned method executes an application. \\
Tool changes during measurement & Measurements are taken only against tagged versions. \\
Study right timing & Study coordinator consulted; extension application filed if advised. \\
\bottomrule
\end{tabularx}
\end{center}

\section{Deliverables}

The thesis; the analyser with two tagged versions corresponding to before and after
attribution; the corpus manifest and the evaluation harness; the generated results
with provenance; a verified bibliography.

\vspace{2em}
\section*{Approval}

\begin{tabular}{@{}p{7cm}p{7cm}}
\\[1.5em]
\rule{6.5cm}{0.4pt} & \rule{6.5cm}{0.4pt} \\
Duy Khanh Tran, student & Prof.\ Petri Vuorimaa, supervisor \\[0.4em]
Date: \rule{3cm}{0.4pt} & Date: \rule{3cm}{0.4pt} \\
\end{tabular}

\end{document}
